\section{Générateurs tritiques, actions quadratiques et transport d’énergie}
\label{sec:hamiltonianos-curvatura}

La géométrie dynamique conserve la distinction entre l’état qui engendre une
réalisation et les coordonnées avec lesquelles celle-ci exprime son évolution. APP
construit et conserve les feuillets additif et multiplicatif ; TRIT oriente le
mode opératoire ; TPK compose sélection, transport et actualisation sur
l’état enrichi au moyen de
\[
 U_t=\operatorname{Upd}_t\circ\operatorname{Tra}_t\circ
 \operatorname{Sel}_t.
\]
L'holonomie nonadique prolonge cet état à travers \(w_6\to w_{12}\to w_{18}\to w_{24}\to w_{30}\to R_{36}\to G_9\). Le retour de phase conserve l'avancement de la mémoire. Ses représentations \(\pi_{\rm HMT},\varphi_{\rm HMT},e_{\rm HMT},\alpha_{\rm HMT}\) appartiennent à la même génération corrélée, avec les lecteurs et l'ordre constructif interne déjà établis. La structure discrète conjointe du continu soutient ensuite les représentations géométriques et physiques. À ce point de la construction, un hamiltonien est la fonctionnelle réalisant un transport déjà typé sur une fibre symplectique ; son domaine, son orientation et son paramètre temporel font partie de l'affirmation.

L’article sur l’extrême et moyenne raison holographique récupère trois
opérateurs concrets. La multiplication par \(4\) modulo \(9\) agit sur
les orbites ternaires d’unités ; dans son module d’augmentation, elle a pour matrice
\[
 C=\begin{pmatrix}0&-1\\1&-1\end{pmatrix},\qquad C^2+C+I=0.
\]
Le transport parabolique et l’incidence aréale--volumétrique sont
représentés, dans leurs plans respectifs, par
\[
 N=\begin{pmatrix}0&1\\0&0\end{pmatrix},\qquad
 F_{\rm av}=\begin{pmatrix}0&1\\1&1\end{pmatrix}.
\]
L'incidence provient du calendrier des feuillets ; ses puissances comptent des compositions de ces modes. Les générateurs réels exprimés sont
\begin{equation}
 J_+=\frac{2C+I}{\sqrt3},\qquad J_0=N,\qquad
 J_-=\frac{2F_{\rm av}^{2}-3I}{\sqrt5},\qquad
 J_\tau^2=-\tau I.
 \label{eq:hc-generadores}
\end{equation}
Ici \(\tau=+1,0,-1\) distingue les régimes elliptique, parabolique et
hyperbolique. Les fibres ont des provenances et des bases propres. Les formules
\(e^{\pi J_+}=-I\), \(e^{J_0}=I+J_0\) et
\(e^{2\log\varphi J_-}=F_{\rm av}^{2}\) évaluent des coordonnées HMT
préalablement engendrées dans ces opérateurs~\cite{hmtX}.

\subsection{Forme alternée et Hamiltonien du générateur}

Dans une fibre réelle orientée de dimension deux, on fixe
\[
 \omega=dQ\wedge dP,\qquad
 \Omega=\begin{pmatrix}0&1\\-1&0\end{pmatrix},\qquad z=(Q,P)^T.
\]
On adopte \(\iota_{X_H}\omega=dH\) ; par conséquent \(\dot Q=\partial_PH\), \(\dot P=-\partial_QH\). Cette convention fixe également le signe de l'action orientée.

\begin{proposition}[Réalisation hamiltonienne d’un générateur de trace nulle]
\label{prop:hc-cuadratica}
Pour une matrice réelle \(J\) de taille deux et de trace nulle, la matrice
\(G_J=-\Omega J\) est symétrique et le Hamiltonien
\begin{equation}
 H_J(z)=\frac12z^TG_Jz
 \label{eq:hc-hj}
\end{equation}
produit exactement \(\dot z=Jz\). Si \(J^2=-\tau I\),
alors \(\det G_J=\tau\). Pour un isomorphisme symplectique
\(B:V\to V'\), le transport
\[
 J'=BJB^{-1},\qquad H'(z')=H_J(B^{-1}z')
\]
satisfait \(G'=B^{-T}G_JB^{-1}\) et entrelace les deux flux.
\end{proposition}
\begin{proof}
Escribiendo \(J=\bigl(\begin{smallmatrix}a&b\\c&-a\end{smallmatrix}\bigr)\), on obtient \(G_J=\bigl(\begin{smallmatrix}-c&a\\a&b\end{smallmatrix}\bigr)\). Ainsi \(\Omega\nabla H_J=\Omega(-\Omega J)z=Jz\). Cayley--Hamilton donne \(J^2=-\det(J)I\), et \(\det(-\Omega)=1\). Enfin, \(B^T\Omega B=\Omega\) permet de substituer \(z=B^{-1}z'\) aussi bien dans la forme que dans l'équation.
\end{proof}

Appliquée à \eqref{eq:hc-generadores}, la proposition donne
\begin{equation}
 H_{J_+}=-\frac{Q^2-QP+P^2}{\sqrt3},\qquad
 H_{J_0}=\frac{P^2}{2},\qquad
 H_{J_-}=\frac{-Q^2-QP+P^2}{\sqrt5}.
 \label{eq:hc-hamiltonianos-nativos}
\end{equation}
Le signe elliptique enregistre l’orientation de \(C\) dans la base d’augmentation.
Le flux inverse \(-J_+\) possède la forme positive opposée. L’ellipse
d’action du corpus utilise précisément une orientation de phase
décroissante pour son Hamiltonien positif. Conserver ce signe rend
l’opérateur cyclique compatible avec l’action \(P\,dQ\).

La conjugaison fournit également un critère exact de comparaison.
Si \(BJ_\tau=J_{\tau'}B\) et \(B\) est inversible, l’élévation au carré
donne \((\tau-\tau')B=0\), donc \(\tau=\tau'\). Un lien peut
conserver \(\omega\) entre fibres de régimes distincts et conserver
leurs mémoires respectives ; la conjugaison des générateurs exige en outre
l’égalité du régime. Cette différence permet d’étudier une transition
tritique sans la remplacer par un changement de nom du même flot.

\subsection{Déformation de l'ellipse et transformation de Legendre}

La coordonnée \(\eta=\operatorname{artanh}(C^*/A)\), dans le domaine \(|C^*/A|<1\), définit le repère symplectique \(M_\eta=\operatorname{diag}(e^{\eta/2},e^{-\eta/2})\). La déformation des deux directions est réciproque et conserve leur produit. Dans la coordinatisation normalisée \((q,p)=M_\eta^{-1}(Q,P)\), considérons
\[
 L_\tau=\begin{pmatrix}0&1\\-\tau&0\end{pmatrix},\qquad
 I_{\tau,\eta}=\frac12\bigl(\tau e^{-\eta}Q^2+e^\eta P^2\bigr).
\]
Ces représentants classent des flux quadratiques orientés. Leur emploi
dans une fibre native s’effectue au moyen d’un repère symplectique déclaré,
conformément à la proposition~\ref{prop:hc-cuadratica}. Le représentant
elliptique positif est celui utilisé par l’ellipse d’action ; pour lui, \(I\)
a la dimension d’une action.

Les changements de repère peuvent s'écrire explicitement. Si \(b_+=\sqrt{\sqrt3/2}\) et \(b_-=\sqrt{\sqrt5/2}\), les matrices
\[
 B_+=\begin{pmatrix}b_+&b_+/\sqrt3\\0&2b_+/\sqrt3\end{pmatrix},
 \qquad
 B_-=\begin{pmatrix}b_-&-b_-/\sqrt5\\0&2b_-/\sqrt5\end{pmatrix}
\]
ont un déterminant égal à un et satisfont \( (-J_+)B_+=B_+L_{+1}\) et \(J_-B_-=B_-L_{-1}\). Pour le régime parabolique, on prend \(B_0=I\). La multiplication des deux colonnes vérifie les égalités, et le déterminant égal à un prouve la conservation de la forme alternée.

\begin{theorem}[Transport variable et lagrangien régulier]
\label{thm:hc-legendre}
Soient \(\eta(t)\) de classe \(C^1\) et \(\omega(t)>0\). L’hamiltonien
\begin{equation}
 H_{\tau,\rm tr}
 =\omega I_{\tau,\eta}+\frac{\dot\eta}{2}QP
 \label{eq:hc-htr}
\end{equation}
conserve \(I_{\tau,\eta(t)}\) le long de ses solutions. Sa
transformation de Legendre par rapport à \(P\) est inversible et donne
\begin{equation}
 \mathcal L_{\tau,\rm tr}(Q,\dot Q,t)
 =\frac{(\dot Q-\dot\eta Q/2)^2}{2\omega e^\eta}
 -\frac{\omega\tau e^{-\eta}}2Q^2.
 \label{eq:hc-ltr}
\end{equation}
\end{theorem}
\begin{proof}
Les équations sont
\[
 \dot Q=\omega e^\eta P+\frac{\dot\eta}2Q,
 \qquad
 \dot P=-\omega\tau e^{-\eta}Q-\frac{\dot\eta}2P.
\]
Dériver \(I_{\tau,\eta(t)}\) et les substituer annule séparément
les termes \(\omega QP\) et les termes de variation de \(\eta\).
La dérivée seconde \(\partial_P^2H=\omega e^\eta>0\)
permet de résoudre
\(P=(\dot Q-\dot\eta Q/2)/(\omega e^\eta)\).
La substitution dans \(P\dot Q-H\) prouve \eqref{eq:hc-ltr}.
De plus,
\[
 P\dot Q-H_{\tau,\rm tr}
 =p\dot q-\frac\omega2(\tau q^2+p^2),
\]
ce qui prouve que le terme mixte est la connexion variationnelle du repère
mobile. La conservation provient du transport complet et non du maintien
des demi-axes constants.
\end{proof}

Le résultat elliptique et son action sont développés dans le corpus variationnel~\cite{hmtX} ; la même preuve exhibe ses deux réalisations quadratiques conjuguées. La nullité du déterminant de la forme parabolique coexiste avec une transformation de Legendre régulière : celle-ci dépend de la dérivée seconde par rapport au moment, et non du déterminant de la forme complète. En revanche, pour le relèvement cotangent linéaire \(H_A(q,p)=p^TAq\), on a \(\partial_p^2H_A=0\). Son action du premier ordre conserve \(\dot q=Aq\) au moyen d'un multiplicateur ; c'est un objet variationnel distinct de \eqref{eq:hc-ltr}.

\subsection{Retour de mémoire et charge variationnelle conservée}

Dans le registre de mémoire du calendrier, les compteurs \((g,s)\)
reçoivent l’accroissement \((4,72)\) au retour complet. La relation
\(72=18\cdot4\) fixe
\[
 v=(1,18)^T,\qquad I_{\rm mem}=s-18g,
 \qquad B=\begin{pmatrix}1&0\\18&1\end{pmatrix}.
\]
Le stabilisateur natif est
\begin{equation}
 N_v=BNB^{-1}=
 \begin{pmatrix}-18&1\\-324&18\end{pmatrix},\qquad
 N_vm=vI_{\rm mem}(m),\qquad N_v^2=0.
 \label{eq:hc-nv}
\end{equation}
Dans le plan réalifié \(m=(g,s)^T\), muni de \(dg\wedge ds\), l'hamiltonien \(H_{\rm mem}=I_{\rm mem}^2/2\) produit \(dm/du=N_vm\). Le paramètre \(u\) appartient à cette réalisation adimensionnelle du registre. La coordinatisation \(q=g,\ p=s-18g\) est symplectique et transforme la un-forme d'action selon
\begin{equation}
 s\,dg=p\,dq+d(9q^2).
 \label{eq:hc-memoria-borde}
\end{equation}
Par conséquent, l’élimination de \(s\) dans l’action du premier ordre donne
\begin{equation}
 \mathcal L_{\rm mem}
 =\frac12\left(\frac{dg}{du}\right)^2
 +18g\frac{dg}{du}
 =\frac12\left(\frac{dg}{du}\right)^2
 +\frac{d}{du}(9g^2).
 \label{eq:hc-lmem}
\end{equation}
L'équation variationnelle est \(d^2g/du^2=0\), et le moment conservé dans la coordinatisation symplectique est \(p=I_{\rm mem}\). Le retour affine complet provient de l'hamiltonien linéaire \(H_{\rm ret}=4I_{\rm mem}\) : son application au temps \(u=1\) est \((g,s)\mapsto(g+4,s+72)\). Les deux hamiltoniens commutent parce qu'ils sont des fonctions de la même charge, mais réalisent des opérations distinctes : stabilisation parabolique et translation du registre. Ces formules restituent le mécanisme variationnel de la mémoire, y compris son terme de frontière~\cite{hmtX}.

\subsection{Courbure métrique et énergie géodésique}

La réalisation métrique des régimes utilise une coordinatisation radiale et une échelle inverse de longueur \(\rho>0\) :
\begin{equation}
 g_{\tau,\rho}=dr^2+S_{\tau,\rho}(r)^2d\theta^2,
 \qquad
 S_{\tau,\rho}(r)=
 \begin{cases}
 \sin(\rho r)/\rho,&\tau=+1,\\
 r,&\tau=0,\\
 \sinh(\rho r)/\rho,&\tau=-1.
 \end{cases}
 \label{eq:hc-metrique}
\end{equation}
La courbure gaussienne est \(K_G=-S''/S=\tau\rho^2\). La coordinatisation est considérée dans son intérieur régulier, \(S>0\) ; pour la réalisation sphérique, \(0<r<\pi/\rho\). Les pôles sont traités avec des coordinatisations régulières différentes. Le paramètre \(\rho\) et la réalisation dimensionnelle de la métrique conservent leur déclaration propre.

Pour une masse réalisée \(m>0\), le lagrangien géodésique et son
Hamiltonien sont
\begin{equation}
 \mathcal L_{\rm geo}=\frac m2(\dot r^2+S^2\dot\theta^2),\qquad
 H_{\rm geo}=\frac1{2m}\left(p_r^2+\frac{p_\theta^2}{S^2}\right).
 \label{eq:hc-geodesique}
\end{equation}
En effet, \(p_r=m\dot r\), \(p_\theta=mS^2\dot\theta\) ;
le Hessien des vitesses est \(m\operatorname{diag}(1,S^2)\)
et est positif et inversible. L’équation radiale est
\(\ddot r-SS'\dot\theta^2=0\), tandis que
\(d(mS^2\dot\theta)/dt=0\) conserve le moment cinétique.
La positivité de cette énergie vaut pour les trois courbures. La forme
indéfinie du générateur hyperbolique de \eqref{eq:hc-hamiltonianos-nativos}
et l’énergie géodésique positive d’une métrique de courbure négative
sont des fonctions sur des espaces distincts. Cette comparaison fixe exactement
le sens de la courbure qui intervient dans chaque assertion.

\subsection{Réponse constitutive et réalisation modale de Maxwell}

Les sections d’action et les coordonnées angulaires précèdent la
réponse constitutive \cite{hmtIntegral}. Dans l’orientation réelle fixée \(x>y>0\),
avec les projecteurs des deux feuillets,
\[
 T=q_+P_++q_-P_-,\quad q_\pm=e^{-x\mp y},\quad
 \mathcal V=(I-T^{90})(I-T^{120})^{-1},
 \qquad 0<q_+<q_-<1,
\]
le calcul fonctionnel produit
\[
 r_\pm=\frac{1-q_\pm^{90}}{1-q_\pm^{120}},\qquad
 \widehat\varepsilon=r_+^2,\quad \widehat\mu=r_-^2,
 \quad \widehat c=(r_+r_-)^{-1},\quad \widehat Z=r_-/r_+.
\]
Les coordinatisations dimensionnelles correspondantes conservent exactement \(\varepsilon\mu=c^{-2}\) et \(\mu/\varepsilon=Z^2\). La permittivité, la perméabilité et la vitesse sont ainsi des représentations du même opérateur, avec les secteurs d'incidence fixés par le TPK ; la dynamique suivante les emploie comme des résultats déjà construits.

Pour exhiber le transport vers un hamiltonien de champ, on fixe une
réalisation spatiale orientée, des conditions de frontière permettant
l’intégration par parties, et un mode transversal réel \(f\) normalisé par
\[
 \nabla\cdot f=0,\qquad \int|f|^2=1,\qquad
 \nabla\times\nabla\times f=k^2f,\qquad k>0.
\]
En jauge transverse et sans sources, \(A=af\) et le moment \(\Pi=pf\) définissent une immersion du plan modal dans l'espace des champs. Pour des constantes constitutives homogènes, l'hamiltonien et la forme symplectique de Maxwell se restreignent à
\begin{equation}
 H_{\rm EM}(a,p)=\frac{p^2}{2\varepsilon}
 +\frac{k^2a^2}{2\mu},\qquad \omega_{\rm EM}=da\wedge dp.
 \label{eq:hc-maxwell-mode}
\end{equation}
L’intégration par parties donne
\(\int|\nabla\times f|^2=k^2\), qui prouve cette restriction.
Avec \(\omega_k=ck\), le changement
\begin{equation}
 Q=\sqrt{\varepsilon\omega_k}\,a,
 \qquad P=\frac{p}{\sqrt{\varepsilon\omega_k}}
 \label{eq:hc-maxwell-symplectic}
\end{equation}
préserve \(da\wedge dp=dQ\wedge dP\) et donne
\(H_{\rm EM}=\omega_k(Q^2+P^2)/2\). Son inverse et ses équations
\(\dot Q=\omega_kP,\ \dot P=-\omega_kQ\)
constituent l’entrelacement modal avec le régime elliptique positif.
La restriction est invariante car l’opérateur de Maxwell conserve le
mode propre transversal. La transformation de Legendre produit
\(\mathcal L_{\rm EM}=\varepsilon\dot a^2/2-k^2a^2/(2\mu)\).

Deux modes transversaux de même fréquence, combinés en quadrature,
réalisent les deux orientations circulaires du corpus. Le champ vérifie
\(\mathbf B=c^{-1}\widehat{\mathbf k}\times\mathbf E\) et
\(\omega=ck\). La relation avec l’oscillateur est démontrée ici par
l’immersion des champs, la normalisation modale et le transport de la forme
symplectique, en maintenant explicites les conditions de frontière.
Un changement de \(k\), de la géométrie de la cavité ou de la préparation
modifie le mode réalisé ; il laisse intacte la généalogie de
\(\varepsilon,\mu,c,Z\).

Dans la coordinatisation lorentzienne du même secteur, le couplage à une section matérielle s'écrit avec \(F=dA\) et une charge \(q_e\) déjà typée. Pour cette réalisation variationnelle, on prend \(\nabla\) égal à la connexion de Levi--Civita de \(g\). Sur \(T^*\mathcal M\), l'hamiltonien covariant
\[
 H_A=\frac1{2m}g^{\mu\nu}(p_\mu-q_eA_\mu)(p_\nu-q_eA_\nu)
\]
est la transformation de Legendre de
\(\mathcal L_A=\frac m2g_{\mu\nu}\dot x^\mu\dot x^\nu
+q_eA_\mu\dot x^\mu\). Le hessien \(mg\) est inversible.
L’équation d’Euler--Lagrange donne
\(m\nabla_u u^\mu=q_eF^\mu{}_{\nu}u^\nu\).
L’antisymétrie de \(F\) conserve \(g(u,u)\). Il s’agit de la
réalisation covariante à paramètre affine, tandis que
\eqref{eq:hc-maxwell-mode} utilise le temps d’une décomposition
spatiale transversale : les deux paramètres et les deux domaines sont définis.

\subsection{Canaux gradués, représentation de Clifford et opérateur de masse}

L'atlas matériel produit d'abord des canaux complets de chemin, de feuillet, de saveur, de couleur et de représentation. À une profondeur finie, sa réalisation hermitienne est \(\mathcal H_K=\bigoplus_{s\in S_K}V_s\), avec une base orthonormale de canaux et une graduation de parité transportée. L'algèbre extérieure du secteur impair engendre les CAR : l'insertion d'un mode \(a_j^*\) et son retrait adjoint \(a_j\) satisfont \(\{a_i,a_j^*\}=\delta_{ij}I\) et \(\{a_i,a_j\}=0\). La preuve compte les signes d'échange de modes complets ; les masses sont évaluées ensuite~\cite{hmtVI}.

Dans le secteur extérieur à un seul mode, la base \((|0\rangle,|1\rangle)\)
identifie l’opérateur de retrait à
\(a=N\). L’application \(e_1\mapsto|0\rangle,\ e_2\mapsto|1\rangle\)
est une isométrie du plan complexe normalisé et entrelace les deux
nilpotents. Cette identification utilise le produit scalaire et la
graduation déclarés ; elle conserve le sens distinct de leurs bases.
Avec le \(i\) interne déjà réalisé, on obtient
\begin{equation}
 \sigma_1=N+N^*,\qquad
 \sigma_2=-i(N-N^*),\qquad
 \sigma_3=NN^*-N^*N.
 \label{eq:hc-pauli}
\end{equation}
La multiplication directe donne
\(\sigma_j^*=\sigma_j\) et
\(\sigma_i\sigma_j=\delta_{ij}I+i\epsilon_{ijk}\sigma_k\).
En particulier, le couple nilpotent et son adjoint produisent une
représentation de Clifford ; l’adjoint est celui de la fibre hermitienne et
est transporté avec elle.

L'opérateur de masse sur une collection finie de canaux \(\mathcal F\) est construit par le registre, le caractère et ses raffinements :
\begin{equation}
 M_\beta=S M_0 S,\qquad
 S=R_{\rm act}^{B_M/2},\quad
 R_{\rm act}=\hbar_{\rm pre}/\hbar_{\rm ret}>0,
 \qquad M_0>0.
 \label{eq:hc-mass}
\end{equation}
Ici, \(B_M\) est l’opérateur diagonal de l’exposant de route et \(M_0\)
contient le caractère raffiné relatif à l’origine électronique.
La positivité découle de \(\langle v,M_\beta v\rangle
=\langle Sv,M_0Sv\rangle>0\) pour \(v\ne0\).
Le produit de feuillets conserve le secteur nul par sa réalisation propre ;
l’exponentielle du caractère matériel reste positive.

Dans cette réalisation, on écrit \(\hbar:=\hbar_{\rm ret}\), la section d'action déjà construite, et \(c:=c_{\rm vac}\) dans la coordinatisation dimensionnelle déclarée. Sur \(\mathbb C^4\otimes\mathcal F\), définissons
\[
 \alpha_j^{D}=\sigma_1\otimes\sigma_j,\qquad
 \beta^D=\sigma_3\otimes I_2.
\]
L'exposant distingue ces matrices de la constante de structure fine. La coordinatisation spatiale \(1+3\) attribue trois composantes d'impulsion \(\mathbf p\). L'application de multiplication
\begin{equation}
 H_D(\mathbf p)
 =c\sum_{j=1}^3\alpha_j^Dp_j\otimes I_{\mathcal F}
 +c^2\beta^D\otimes M_\beta
 \label{eq:hc-dirac}
\end{equation}
est la réalisation libre de Dirac de cet opérateur matériel.

\begin{theorem}[Factorisation spinorielle de la dispersion matérielle]
\label{thm:hc-dirac-square}
La matrice \eqref{eq:hc-dirac} est autoadjointe et satisfait
\begin{equation}
 H_D(\mathbf p)^2
 =I_4\otimes\bigl(c^2|\mathbf p|^2I_{\mathcal F}
 +c^4M_\beta^2\bigr).
 \label{eq:hc-dirac-square}
\end{equation}
Dans un sous-espace propre \(M_\beta v=m v\), ses énergies sont \(\pm E_m(\mathbf p)\), où \(E_m=\sqrt{c^2|\mathbf p|^2+m^2c^4}\).
\end{theorem}
\begin{proof}
Les relations de \eqref{eq:hc-pauli} donnent
\(\{\alpha_i^D,\alpha_j^D\}=2\delta_{ij}I_4\),
\(\{\alpha_i^D,\beta^D\}=0\) et \((\beta^D)^2=I_4\).
Le facteur de saveur commute avec les matrices spinorielles.
Les termes croisés s’annulent et les carrés donnent
\eqref{eq:hc-dirac-square}. La trace nulle et les relations de Clifford
produisent les deux branches, avec une multiplicité spinorielle de deux lorsque \(E_m>0\).
\end{proof}

Dans la représentation des moments, le domaine naturel est
\(\{\psi\in L^2(\mathbb R^3;\mathbb C^4\otimes\mathcal F):
H_D\psi\in L^2\}\). Étant une matrice hermitienne mesurable agissant
par multiplication, l’opérateur est autoadjoint sur ce domaine. La
représentation spatiale provient de la transformation unitaire de Fourier
postérieure à la construction des canaux. Le théorème fixe la réalisation
libre et sa dépendance exacte à la masse HMT ; l’interaction conserve ses
opérateurs de connexion et son domaine propre.

Dans la branche positive \(m>0\), la transformation de Legendre ordinaire du symbole \(E_m(\mathbf p)\) est
\[
 \mathbf v=\frac{c^2\mathbf p}{E_m},\qquad
 \mathcal L_m=-mc^2\sqrt{1-|\mathbf v|^2/c^2},
 \qquad |\mathbf v|<c.
\]
Résoudre \(\mathbf p=m\mathbf v/\sqrt{1-|\mathbf v|^2/c^2}\) et substituer dans \(\mathbf p\cdot\mathbf v-E_m\) démontre l'égalité. L'action spinorielle, en revanche, est du premier ordre dans le champ : \(\mathcal L_D=\frac{i\hbar}{2}(\psi^*\dot\psi-\dot\psi^*\psi)
-\psi^*H_D\psi\). Sa variation donne \(i\hbar\dot\psi=H_D\psi\) ; son hessien par rapport aux vitesses du champ est singulier. Les deux affirmations sont compatibles car elles agissent sur des espaces variationnels différents.

Pour une connexion abélienne statique dans une coordinatisation spatiale plane, avec \(\Pi_j=-i\hbar\partial_j-q_eA_j\), une masse constante et un potentiel scalaire nul, on conserve sur les fonctions lisses à support compact
\[
 [\Pi_i,\Pi_j]=iq_e\hbar\epsilon_{ijk}B_k,
\]
et la même factorisation produit
\begin{equation}
 H_D(A)^2=c^2\boldsymbol\Pi^2+c^4M_\beta^2
 -q_e\hbar c^2\boldsymbol\Sigma\cdot\mathbf B,
 \qquad \Sigma_j=I_2\otimes\sigma_j.
 \label{eq:hc-dirac-curvature}
\end{equation}
Les facteurs identité de saveur restent implicites. Le terme de courbure est le produit des deux antisymétriques : le commutateur des dérivées covariantes et la multiplication de Pauli. Son coefficient est fixé lors du choix de la représentation de charge et de la coordinatisation d'action ; il n'est pas sélectionné au moyen d'une valeur spectrale cible.

Si \(V\) est la transformation unitaire de saveur déjà construite, alors
\[
 (I_4\otimes V)H_D(M_\beta)(I_4\otimes V^*)
 =H_D(VM_\beta V^*).
\]
Le mélange unitaire conserve le spectre complet. La congruence positive
de \eqref{eq:hc-mass} transporte une autre structure et peut modifier les
valeurs propres dans une métrique fixe. Par exemple,
\(M_0=\operatorname{diag}(1,2)\), \(S=\operatorname{diag}(2,1)\)
donne \(SM_0S=\operatorname{diag}(4,2)\). Si l’on transporte également le
produit scalaire vers \(G'=S^*S\), le problème généralisé
\(SM_0Sv=\lambda G'v\) retrouve les valeurs propres de \(M_0\)
au moyen de \(w=Sv\). Ainsi, le changement de base de saveur, la déformation
matérielle et le transport de métrique ont des conséquences spectrales précises.

\subsection{Courbure du potentiel scalaire dans la réalisation électrofaible}

Le secteur électrofaible utilise le déterminant du même opérateur angulaire, \(D_A=\det T=e^{-2A_{\rm rad}}\), avec les signatures matérielles déjà construites. Dans la réalisation angulaire avec raffinements communs de \(W\) et \(Z\), leurs signatures \((94,-1)\) et \((95,-1)\) donnent \(m_W^\angle/m_Z^\angle=e^{-A_{\rm rad}}\). Dans la convention rationalisée à l'ordre de l'arbre du corpus,
\[
 g^2=\frac{4\pi_{\rm HMT}\alpha_{\rm HMT}}{1-D_A},\qquad
 \lambda_H=\frac{\pi_{\rm HMT}\alpha_{\rm HMT}}{2(1-D_A)}
 \exp\!\left(6A_{\rm rad}+\frac{10}3C^*_{\rm rad}\right)>0.
\]
La seconde formule résulte de la signature scalaire \((97,9)\) et de
\(m_H^2/m_W^2=8\lambda_H/g^2\). La normalisation du potentiel est
\(V(\mathsf H)=-\mu_H^2\mathsf H^*\mathsf H
+\lambda_H(\mathsf H^*\mathsf H)^2\). Ces compositions conservent
leur échelle et leur convention d’arbre ; les raffinements qui s’annulent
dans un rapport sont déclarés dans la réalisation elle-même.

Dans une coordinatisation radiale réelle \(\mathsf H^*\mathsf H=h^2/2\), le potentiel est \(V(h)=-\mu_H^2h^2/2+\lambda_Hh^4/4\). Pour \(\mu_H^2>0\), ses points critiques sont \(0\) et \(\pm v\), avec \(v^2=\mu_H^2/\lambda_H\), et
\[
 V''(0)=-\mu_H^2,\qquad V''(\pm v)=2\mu_H^2.
\]
Dans les unités naturelles de cette coordinatisation, le mode homogène normalisé possède \(H_h=p_h^2/2+V(h)\). Sa linéarisation en un point critique \(h_0\) a pour générateur
\[
 A_{h_0}=\begin{pmatrix}0&1\\-V''(h_0)&0\end{pmatrix},
 \qquad A_{h_0}^2=-V''(h_0)I.
\]
Par conséquent, l'origine est hyperbolique et les minima sont elliptiques ; la frontière \(V''=0\) est de type parabolique. Cette affirmation se déduit de la dérivée seconde du potentiel et d'une normalisation modale explicite. L'énergie matérielle positive du minimum et la direction instable d'un autre fond appartiennent ainsi à la même fonction avec des points de développement différents. Le mode linéarisé est un transport particulier des trois types quadratiques ; ses fréquences et ses unités proviennent du potentiel, et non d'une identification de son hessien à la courbure gaussienne de \eqref{eq:hc-metrique}.

\subsection{Mode de courbure de Yang--Mills et oscillation elliptique}

La réalisation de Yang--Mills construite dans les sections précédentes
possède un générateur positif \(D^{\rm YM}\) de dimension inverse de
longueur. L’échelle
\[
 \kappa_{\rm av}=\frac{\mu_{\rm vol}}{\ell_0}
 \left(\frac{\pi_{\rm HMT}}{\varphi_{\rm HMT}^2}-1\right)>0,
 \qquad \delta=3\kappa_{\rm av}
\]
provient de l'incidence des feuillets et du lecteur régional, avec l'unité longitudinale \(\ell_0\) déclarée. Le vide est simple et le témoin de courbure dans la coordinatisation produit est
\[
 W(q)=\mathbf1_{\{q_1+\cdots+q_6=0\}}-\frac13,
 \quad \mathbb EW=0,\quad \|W\|^2=\frac29,
 \quad D^{\rm YM}W=\delta W.
\]
Le transport unitaire vers la mesure d’interaction transporte \(W\) en
\(W_g\) et conserve ces identités. Les inclusions de raffinement
\(J_m\) entrelacent les générateurs et conservent le témoin, de sorte que
celui-ci demeure dans la limite physique. On utilise ici le résultat complet
des preuves du vide, de l’écart uniforme et de la reconstruction, non une
valeur propre isolée d’une matrice finie~\cite{HMT-YM}.

Soit \(w=W_g/\sqrt{2/9}\) le vecteur unitaire dans la fibre physique et
\(\hbar=\hbar_{\rm ret}>0\). Le Hamiltonien énergétique est
\(H_E=\hbar cD^{\rm YM}\). Sur l’espace de Hilbert réalifié,
on adopte la forme \(\omega_{\mathcal H}(u,v)=2\hbar
\operatorname{Im}\langle u,v\rangle\), avec un produit scalaire antilinéaire
par rapport au premier argument. Définissons
\begin{equation}
 \mathfrak F(Q,P)=\frac{Q+iP}{\sqrt{2\hbar}}\,w.
 \label{eq:hc-ym-embedding}
\end{equation}

\begin{theorem}[Immersion symplectique du mode physique de courbure]
\label{thm:hc-ym-oscillator}
L’application \(\mathfrak F:\mathbb R^2\to\mathcal H^{\rm phys}\)
est une immersion réelle dont l’image est invariante par l’évolution physique et
vérifie
\begin{equation}
 \mathfrak F^*\omega_{\mathcal H}=dQ\wedge dP,
 \qquad
 \langle\mathfrak F,H_E\mathfrak F\rangle
 =\frac{c\delta}{2}(Q^2+P^2).
 \label{eq:hc-ym-energy-pullback}
\end{equation}
L'équation de Schrödinger restreinte entrelace exactement le flux elliptique \(\dot Q=c\delta P,\ \dot P=-c\delta Q\). L'application est compatible avec les inclusions de raffinement.
\end{theorem}
\begin{proof}
Pour deux variations \((a,b),(a',b')\), le produit scalaire de leurs
images a pour partie imaginaire \((ab'-ba')/(2\hbar)\).
Cela prouve la première identité. L’équation propre et \(\|w\|=1\)
donnent la seconde. Puisque
\(\dot\psi=-icD^{\rm YM}\psi\), le coefficient \(Q+iP\)
satisfait \(\dot Q+i\dot P=-ic\delta(Q+iP)\).
Enfin, \(J_mw_m=w_{m+1}\) implique
\(J_m\mathfrak F_m=\mathfrak F_{m+1}\). Ces identités passent à la
limite par les isométries déjà construites.
\end{proof}

Le plan complet décrit les amplitudes d’un mode ; son cercle
\(Q^2+P^2=2\hbar\) représente les vecteurs normalisés de cet
espace propre complexe. Sur ce cercle, la phase globale évolue et le
rayon quantique reste fixe. Les superpositions d’espaces propres distincts
conservent, en outre, leurs phases relatives. L’immersion maintient cette
différence entre l’espace des amplitudes et l’espace projectif.

L'énergie associée au saut spectral et le saut de masse sont \(\Delta_E=\hbar c\delta\) et \(m_{\rm gap}=\hbar\delta/c\). La fréquence modale \(c\delta\) a la dimension d'un temps inverse. Les longueurs euclidiennes de \(e^{-sD^{\rm YM}}\) et le temps physique de \(e^{-itH_E/\hbar}\) sont reliés par la coordinatisation et la reconstruction, en conservant leurs groupes et semi-groupes distincts. La composition \(\mathfrak F\circ M_\eta^{-1}\) transporte cette même énergie vers l'ellipse d'action ; si le repère varie, le terme de connexion de \eqref{eq:hc-htr} complète son évolution.

Le lien obtenu est une réalisation exacte sur le mode invariant de
courbure et ses prolongements. La positivité de l’hamiltonien quantique
engendre une rotation dans sa réalification symplectique ; la dénomination
hyperbolique d’une métrique spatiale appartient au tenseur métrique, et le
transport unipotent du registre appartient à sa mémoire. En dimension
finie, un groupe \(I+tN\) unitaire pour tout \(t\in\mathbb R\), avec
\(N^2=0\), impose \(N=0\) : dériver l’unitarité donne \(N^*=-N\),
et alors \(N^*N=-N^2=0\). Le flot parabolique classique conserve sa
forme symplectique et possède sa quantification dans des espaces appropriés ; cette
identité délimite seulement l’identification à une matrice unitaire finie.

Ainsi s’articule un ensemble de transports concrets. La mémoire
produit une charge et une action parabolique ; le repère de l’ellipse produit
la connexion variationnelle et son lagrangien ; les réponses constitutives
déterminent les coefficients de Maxwell ; la graduation des canaux construit
Clifford et permet de factoriser l’opérateur matériel ; la reconstruction de
Yang--Mills apporte une immersion modale symplectique qui conserve exactement
son écart et son raffinement. Chaque identification dispose d’une application,
d’une forme conservée et d’une équation entrelacée.
